% Auto-generated by scripts/06_emit_structure.py
% Source: scans 013–017
% Section id: frontmatter/preface
% Phase 3 deliverable. Footnotes (Phase 4), citations (Phase 5),
% and Wade-Giles → Pinyin (Phase 6) are not yet applied here.

\chapter*{Preface}
\label{fm:preface}

% source: scan 013, printed xiii
It is the first responsibility of historians-though not their last-
t-to assess their
sources and to make their critical methods and conclusions about those sources available
to other scholars. In the case of Shang China, historians must strive to discern a series of
consistent patterns in divination inscriptions which were carved on some one hundred
thousand scapula and plastron fragments over three thousand years ago. The historiographical problems are the standard ones involving questions of restoration, dating, decipherment, chronological pattern, authenticity, and motivation. Technical decisions about
such matters can significantly alter our understanding of Shang reality, for the line between
a study of historical sources and the writing of history using those sources is exceedingly
fine, especially in the field of ancient history where the comparative paucity of documents
gives great significance to each one. We can never fully understand Shang history or any
other as it really was, but we must, as a first step in our study of the Shang, understand the
documents as they really are. It is the purpose of this book to introduce the student to
oracle inscriptions, to define their nature, to show how they may be deciphered, and to
suggest how they may be used as historical sources.
Sources of Shang History is focused on what I call the historical period, the era of the
eight or nine kings from Wu Ting to Ti Hsin (ca. 1200-1050 B.C., in my view) for which
we have contemporary Shang inscriptions.\footnote[1]{On the question of whether or not Lin Hsin reigned, see table 1, note h. On the Shang dates, see appendix 4.}\footnote[2]{On these terms, see Soper (1966) and, for the problems involved, Kane (1973), pp. 356-358.}\footnote[3]{E.g., Chang Kwang-chih (1968), p. 234; cf. KK (1975.1), pp. 31-32.}\footnote[4]{See ch. 4, n. 24.}\footnote[5]{The removal forms the occasion for the ``P'an Keng” chapter of Shang-shu; on its probable date of composition, see Creel (1938), pp. 63-69; Ch'en Meng-chia (1956), p. 251; Matsumoto Masaaki (1966), pp. 420-433; Wheatley (1971), pp. 1314, 420; cf. Ho (1975), p. 307, however, who finds that much of its content is authentic, i.e., Shang. X222}

% source: scan 014, printed xiv
a place called Yin, we are not sure where Yin was located;\footnote[6]{See ch. 2, n. 2. Accepting the Ku-pen chu-shu 7. This theory appears first in Ti-wang shih-chi chi-nien statement that this was the last Shang, a work of the third century A.D., as capital removal (Fan [1962], p. 21), and assuming that the 安陽 site (focused on what is thought to have been the temple-palace complex at Hsiaot'un) was the Shang capital in the historical period, modern scholars have generally assumed that P'an Keng removed the capital there (e.g., Creel [1938], pp. 135-139 and Ch'en Meng-chia [1956], p. 252; Tung [1945], pt. 2, ch. 5, p. 23a, disagreed). But there are possible difficulties with this view; to cite but two of them: Ssu-ma Ch'ien states that P'an Keng moved south of the Yellow River (SKK, ch. 3, p. 20; Chavannes [1895], p. 194), yet 安陽 is to its north and, so far as we know, always has been; further, the claim that there were no more capital removals after P'an Keng is contradicted by other traditions (cited by Ch'en Meng-chia [1956], p. 252; cf. Chou Hung-hsiang [1958], p. 22, and the summary by Lefeuvre [1975], p. 64, n. 1).} 6 (4) the oracle inscriptions
make no reference to any great settlement (ta-yi ) called Yin, though they do refer
to the ta-yi Shang, a fact which undercuts the traditional view that the dynasty was known
as Shang prior to P'an Keng's removal and as Yin after it. Given these uncertainties, I
refer to the period from Wu Ting (not P'an Keng) to Ti Hsin as the historical period.
When I refer to the Shang state I mean the polity ruled by the kings of the historical period
whose capital, or cult center, is represented by the foundations, tombs, and other remains
at 安陽 that are contemporary with the oracle inscriptions excavated there. It may
also be remarked that we are not absolutely certain that the Shang called themselves Shang,
rather than Yin; thus uncertainty extends even to the title of this book. I use the term
Shang because it seems probable that the ta-yi Shang did refer to the capital and thus to the
dynasty.\footnote[8]{This view is supported by the inscription on the Ho-tsun, excavated in 1963, and datable to the fifth year of Ch'eng Wang of Chou, which contains the retrospective phrase, wei Wu Wang chi k'e ta-yi Shang, “it was when Wu Wang had conquered the great settlement Shang'' (T'ang Lan [1976], p. 60). For a general introduction to the problem, see Barnard (1960b), pp. 499-500, who notes that ``Yin” is never applied to the dynasty in scientifically excavated bronzes, and 屈萬里 (1965a), pp. 88-89, who provides a useful tabulation of references to Yin and Shang in Chou texts. Ke Yi-ch'ing (1939) is interesting on the origins of the name Shang itself. U}8
Sources is not a book to be enjoyed at one sitting, but a book to be used and reused by
those who are working with, or want to work with, the oracle-bone inscriptions. As an initial
study, it has few discrete edges, but unravels in numerous footnotes that suggest the need
for further research. I frequently find occasion to cite unpublished doctoral dissertations,
conference papers, or research in progress. I do this both to give credit to the numerous
scholars whose endeavors have aided my own and to indicate the kinds of problems on
which scholars are currently working. My own research in progress deals with such topics
as the status of the diviners, the theology of Shang divination, how the cracks were read,
the reason why the oracle bones were carved and stored, the controversy over the Old
and New Schools, and the absolute date of the dynasty. I refer to these future publications
by the general title Studies in Shang Divination (abbreviated Studies).
It may be helpful to introduce the various conventions of citation and transcription
quoted in T'ai-p'ing yü-lan, ch. 83, p. 8b: ``The
Emperor P'an Keng removed the capital to Yin;
for the first time one changed Shang and said Yin.''
For oracle inscriptions about the ta-yi Shang, see
\$279.2-3.

% source: scan 015, printed xv
used in this book. With the exception of the references to Shima's Inkyo bokuji sōrui
(abbreviated as S) which omit the p., all page references are preceded by p. or pp. Numbers
not preceded by these letters refer to a rubbing number, drawing number, entry number,
and so on; thus, 顏一萍 (1961), pp. 207-215 is not the same as 顏一萍 (1961),
207-217. The method by which oracle-bone inscriptions are cited is explained in ch. 3, n. 5.
The numbers used to identify oracle-bone inscriptions are italicized when the relative date of
the inscription is germane to the discussion and can be assigned with some certainty (see
ch. 4, n. 1). The letter D indicates that the citation is to a drawing of an oracle inscription
rather than to a rubbing or photograph.\footnote[9]{In some cases, rubbings have recently become available for inscriptions which had previously only been published as drawings; rubbings of the inscriptions in Nan-pei, “Ming,'' and certain inscriptions in \pinyinterm{kufang}, for example, may now be found in Ming-hou and US respectively (for other cases, see ch. 5, n. 56). In order to facilitate finding the inscriptions in Shima Kunio's Inkyo bokuji sōrui (sec. 3.3.2), however, it is still the drawings that are cited in this book; the equivalent rubbings are given in the finding list. XV} 9 The symbols used to indicate missing and restored
characters are described in sec. 5.6. Archaic pronunciations, generally Karlgren's, are
preceded by an asterisk (see ch. 3, n. 46).
The footnotes, like the dedication page, reveal the extent of my debt to the large
number of scholars, primarily Chinese and Japanese, who have devoted themselves to
oracle-bone studies. My debt to contemporaries who have given me the benefit of their
criticisms and generously shared their knowledge is equally large. In particular, I should
like to thank Noel Barnard, Chang Kwang-chih, Chou Hung-hsiang, Hsü Chin-hsiung,
J. A. Lefeuvre, Gilbert Mattos, Stanley Mickel, Henry Rosemont, Jr., and Takashima
Kenichi. I am happy to acknowledge my special debt to David S. Nivison. We have tried
``cracking'' many inscriptions together; his sensitive and original readings have taught me
much. None of these scholars is responsible for the opinions and conclusions given in the
book. I should like to thank Jack Jacoby of Columbia University's East Asiatic Library,
Barbara Stephen and Hsü Chin-hsiung of the Royal Ontario Museum, Vadime Elisséeff
of the Musée Cernuschi, Michèle Pirazolli-t'Serstevens of the Musée Guimet, Françoise
Wang, the librarian of the Institut des hautes études chinoises of the Collège de France,
Jean Schmitt of the Metropolitan Museum of Art, and Don W. Dragoo of the Carnegie
Museum of Natural History, for allowing me to examine the inscribed oracle-bone fragments
in their collections. I owe a debt of thanks to James F. Berry of the Department of Biology
of the University of Utah for making the turtle identifications and for writing appendix 1, to
John M. Legler of that same department for providing the facilities for Berry's research,
and to Richard Kunst and Douglas McOmber, who made the plastron measurements.
Graduate students have contributed much to my understanding of the Shang, both as
seminar participants and as research assistants. I should like to thank Dessa Bucksbaum,
Sue Currie, Inge Dietrich, Sue Glover, Morgan Jones, David Kamen, David Keegan,
Richard Kunst, Bonnie Ng, John W. Olsen, and Beth Upton. I am grateful to David
Keegan, Richard Shek, and Ronald P. Toby for helping with Bibliography B, and to Dessa

% source: scan 016, printed xvi
Bucksbaum for preparing the finding list and index. A more intangible academic debt is
owed to Chao Lin, formerly of the Academia Sinica, who aroused my interest in the Shang
inscriptions and guided my first faltering steps. And my awareness of the importance of
historiographical problems, the rigor with which they should be pursued, and the rewards
of doing so owes much to Hans Bielenstein and C. Martin Wilbur, my teachers at Columbia
University. It is a special pleasure to thank my many colleagues and friends at Berkeley-in
particular, Delmer Brown, Irwin Scheiner, Frederic Wakeman, and the other Far Eastern
historians, Eric Gruen and Raphael Sealey in Classical History, John Cikoski in Oriental
Languages, and Wolfram Eberhard in Sociology-whose support and critical interest
have encouraged and guided my efforts.
The staff of the East Asiatic Library at Berkeley has been an unfailing source of
informed assistance. The Library Delivery Service of the Main Library at Berkeley also
deserves special thanks for its speed and courtesy. I am grateful to Shima Kunio and his
publishers for permission to reproduce transcriptions from his Inkyo Bokuji Sōrui.
Finally, I should like to express my gratitude to the American Council of Learned
Societies, to the Committee on Research and to the Humanities Institute of the University
of California at Berkeley, and to the National Endowment for the Humanities for providing
grants which enabled me to devote my time to the writing of this book.
D.N.K.
Berkeley, 12 January 1977
Well over a year and a half having passed since it became impossible to make significant
changes in the manuscript, I should like to refer the reader to recent developments, archaeological and scholarly, which bear on issues discussed in this book.
5,
The excavation at 小屯 in the spring of 1976 of tomb number undisturbed
and lavishly furnished with grave goods, was an unprecedented archaeological event (KK
1977.3; 1977.5; a further report is due to appear in KKHP 1977.2). (Bibliographic references
to the postscript are given on pp. 255-256.) The name of Fu Hao appears on over
sixty of the bronze vessels, and the temple name, Ssu Mu Hsin (to be read perhaps
as Hou [?] Mu Hsin or Kou [?] Hsin) appears on four. These inscribed names not
only identify, it is presumed, the main occupant of the tomb, but also bear on the question
of when the Shang first started to cast inscribed bronze vessels (p. 134, n. 2, below). Most
scholars believe the tomb to the late period I burial of the Fu Hao who was Wu Ting's
consort (see the oracle-bone inscriptions at S139.4-141.1). This conclusion, whose final
acceptance must await the full report of the grave and its contents, bears in turn on the
absolute date of Wu Ting's reign. If, as Cheng Chen-hsiang and Shih Chih-lien

% source: scan 017, printed xvii
suggest (KK 1977.5, pp. 341, 348), the tomb of Fu Hao was roughly contemporary
with the great tomb at Wu-kuan-ts'un (WKGM 1), so that it too may be dated to late
period I, rather than to period II or later, then the corrected carbon-14 date of 1241 ±
B.C. for WKGM 1 (table 2, sample ZK-5) accords well with the date of 1200-1181 B.C.
assigned to period I by the revised short chronology of table 38 (on these points see p. 173,
n. 13, and p. 176, n. 28). (Scholars interested in carbon-14 dates, including those given in
table 2, may now consult the convenient summary provided by Hsia Nai [1977]; his
discussion bears on the relative age of Neolithic sites at which pyromantic remains have
been found [p. 5, n. 14, below].)
More recent, and equally extraordinary, was the discovery, some 100 kilometers
west of Sian, of over 15,000 oracle-bone fragments (mainly turtle-shell), supposedly divined
by the Western Chou. According to the initial news report, 127 of the shell fragments bore
writing. The inscriptions, which are said to record that Wen Wang of Chou sacrificed to
Ti Yi of Shang, and that the Shang king came to the Chou homeland in Shensi (Hsinhua-she, 1977), will, when published, presumably throw light on Shang-Chou relations (cf.
p. 180, n. 23); if the inscriptions contain the term Yin-wang E, as reported, then the Chou
appear to have referred to the Shang as Yin prior to the conquest (cf. p. xiv, n. 8). The find
also indicates that the Late Shang practice of cracking and inscribing large numbers of
oracle bones may have been more widespread than the present corpus of inscriptions,
excavated primarily in the 小屯 area, would suggest. (cf. my comments on p. 138).
It should be noted that Chang Kwang-chih (1977) has superseded Chang (1968) as
the basic introduction in English to the field of Shang archaeology and culture.
In the realm of oracle-bone scholarship, what I have called the t'ou period (p. 43, n. 79)
is the subject of a new study by Lung Yü-ch'un (1976). The function of the word represented
by the graph yu + (p. 78, n. 85) has been further discussed by Nivison (1977a), whose
conclusions (first advanced in Nivison [1977]) have in turn been challenged by Takashima
(1977a). Yü Hao-liang (1977) argues that the crack notation usually read as hung-chi
(p. 121) should be read as yin-chi 1, ``prolonged auspiciousness.” The publication of
Mickel (1977a) provides indexes to rejoined and reduplicated fragments in Britton (1935),
Creel (1937), \pinyinfull{kufang}, \pinyinfull{baiken}, \pinyinfull{cuibian},
\pinyinfull{wenlu}, and the three \pinyinfull{yezhong-pianyu} collections
(see \pinyinfull{yezhong-pianyu-chuji} and \pinyinfull{yezhong-pianyu-sanji} in Bibliography A);
this supplements his other indexes cited at p. 150,
n. 66, and in table 9. Chou Hung-hsiang's initial use of the computer to match oracle-bone
fragments (p. 150, n. 68) has been carried a step further by scholars in Peking (Tung,
Chang, and Ch'en [1977]). Chiang Hung (1976), dealing with the question of whether
recently discovered Shang sites in the middle Yangtze region were part of the Shang state,
Hu Hou-hsüan (1977), dealing with the subject of bird totemism, and Shen Wen-cho (1977),
dealing with slaves and their agricultural role, should be added to the list of recent works
(p. 154, n. 87) that use the inscriptions to study Shang history.
January 1978
